\documentclass[11pt,a4paper]{article}
\usepackage[T1]{fontenc}
\usepackage{lmodern}
\usepackage[margin=28mm]{geometry}
\usepackage{amsmath,amssymb,amsthm,mathtools,mathrsfs}
\usepackage{microtype,booktabs,array,enumitem}
\usepackage{xcolor}
\definecolor{darkblue}{rgb}{0,0,0.7}
\newcommand{\defn}[1]{\textsl{\color{darkblue} #1}}
\usepackage[colorlinks=true,linkcolor=darkblue,citecolor=darkblue,urlcolor=darkblue]{hyperref}
\hypersetup{pdftitle={Equivariant magnitude, Borel constructions, and distance algebras},pdfauthor={}}
\numberwithin{equation}{section}
\newtheorem{theorem}{Theorem}[section]
\newtheorem{proposition}[theorem]{Proposition}
\newtheorem{corollary}[theorem]{Corollary}
\newtheorem{lemma}[theorem]{Lemma}
\theoremstyle{definition}
\newtheorem{definition}[theorem]{Definition}
\newtheorem{example}[theorem]{Example}
\newtheorem{convention}[theorem]{Convention}
\theoremstyle{remark}
\newtheorem{remark}[theorem]{Remark}
\newcommand{\Z}{\mathbb Z}
\newcommand{\Q}{\mathbb Q}
\newcommand{\C}{\mathbb C}
\newcommand{\F}{\mathbb F}
\newcommand{\G}{G}
\newcommand{\MH}{\operatorname{MH}}
\newcommand{\MC}{\operatorname{MC}}
\newcommand{\Mag}{\operatorname{Mag}}
\newcommand{\Tor}{\operatorname{Tor}}
\newcommand{\Ext}{\operatorname{Ext}}
\newcommand{\Hom}{\operatorname{Hom}}
\newcommand{\End}{\operatorname{End}}
\newcommand{\RHom}{\mathbf{R}\!\operatorname{Hom}}
\newcommand{\Tot}{\operatorname{Tot}}
\newcommand{\Rep}{\operatorname{Rep}}
\newcommand{\Irr}{\operatorname{Irr}}
\newcommand{\Res}{\operatorname{Res}}
\newcommand{\Ind}{\operatorname{Ind}}
\newcommand{\tr}{\operatorname{tr}}
\newcommand{\rad}{\operatorname{rad}}
\newcommand{\Nat}{\operatorname{Nat}}
\newcommand{\im}{\operatorname{im}}
\newcommand{\id}{\operatorname{id}}
\newcommand{\derivedtensor}[1]{\mathbin{\mathop{\otimes}\limits^{\mathbf L}_{#1}}}
\newcommand{\Bor}{\mathrm{Bor}}
\newcommand{\Br}{\mathrm{Br}}
\newcommand{\RC}{R_{\C}}
\newcommand{\Burn}{\mathcal A}
\newcommand{\Orb}{\mathcal O}
\newcommand{\one}{\mathbf 1}
\newcommand{\coeff}{\subsection*{\color{darkblue}Coefficients}}
\newcommand{\source}{\subsection*{\color{darkblue}Source and status}}
\setlist{itemsep=2pt,topsep=5pt}
\setlength{\emergencystretch}{2em}

\title{Equivariant magnitude, Borel constructions,\texorpdfstring{\\}{ }and distance algebras}
\author{}
\date{Research note --- 15 September 2026}

\begin{document}
\maketitle
\begin{abstract}
We compare the representation-valued equivariant refinement of magnitude
(co)homology, its Borel version, and their realizations over skew group
algebras of distance algebras.  We then construct an integral
Burnside-ring-valued magnitude series and describe a possible refinement
using Bredon (co)homology.  Ground-ring assumptions, module structures,
and the difference between fixed points and invariant vectors are kept
explicit.  Ordinary magnitude (co)homology and the Asao--Ivanov algebraic
model are assumed throughout.  Results taken from the literature are
identified at their point of use; equivariant specializations and
additional deductions are proved here.  No claim of novelty is intended.
\end{abstract}
\tableofcontents

\section{The usual representation-valued equivariant theory}
\label{sec:ordinary}

\subsection{Standing conventions and coefficients}

Let $X$ be a finite metric space, and let $\G$ be a finite group acting
on $X$ by isometries.  Graphs are treated through their vertex metrics;
the letter $\G$ always denotes the symmetry group.  Unless explicitly
stated otherwise, all distances are finite.  The empty metric space is
allowed as a fixed-point space, and its magnitude and all its magnitude
groups are zero.  Every fixed subset carries the \defn{restricted metric}.
In particular, for a graph, this need not be the shortest-path metric of
the induced subgraph on its fixed vertices.

Let $\Lambda\subseteq\mathbb R_{\geq0}$ be the additive monoid generated
by the positive distances in $X$, and let $D\subseteq\mathbb R$ be the
additive group they generate.  For a coefficient ring $K$, write
$K[[q^\Lambda]]$ for formal sums $\sum_{\ell\in\Lambda}a_\ell q^\ell$,
with multiplication $q^a q^b=q^{a+b}$.  There are only finitely many
elements of $\Lambda$ in any bounded interval, so convolution is
well-defined.  If the metric is integral, this is naturally a subring
of $K[[q]]$.  All series below are formal; no analytic substitution for
$q$ is required.  When comparing several spaces we use the monoid
generated by all their distances.

For a commutative unital ring $R$, use the notation
\[
 C_{*,\ell}(X;R)=\MC_{*,\ell}(X;R),\qquad
 C^{*,\ell}(X;R)=\Hom_R(C_{*,\ell}(X;R),R).
\]
We assume the ordinary definitions and products, as in
\cite{HW,LS,Hepworth,AI}, without recalling them.  We need only name the
standard basis set
\begin{equation}\label{eq:tuple-set}
 \mathcal T_{n,\ell}(X)=
 \left\{(x_0,\ldots,x_n): x_{i-1}\ne x_i,\quad
             \sum_{i=1}^n d(x_{i-1},x_i)=\ell\right\}.
\end{equation}
Thus $C_{n,\ell}(X;R)=R[\mathcal T_{n,\ell}(X)]$.
If $\delta>0$ is the least positive distance, then
$\mathcal T_{n,\ell}(X)=\varnothing$ for $n>\ell/\delta$.
The cases $|X|\leq1$ satisfy the same boundedness assertion directly.
Consequently every fixed-length chain complex is bounded and finite free
over $R$.  None of this requires a field.

We use $\Mag(X;q)$ for ordinary formal magnitude.  Its graded Euler
characteristic description is the one in \cite[Theorem 8]{HW} for graphs
and \cite[Corollary 7.15]{LS} for finite metric spaces.  In particular,
it is the alternating series of the cardinalities in
\eqref{eq:tuple-set}, with integral coefficients.  Over $\Z$, passing
to ranks of homology has the same effect as tensoring homology with
$\Q$; this numerical invariant does not record torsion.

\subsection{Actions and the representation ring}

The action on tuples is diagonal:
\begin{equation}\label{eq:tuple-action}
 g(x_0,\ldots,x_n)=(gx_0,\ldots,gx_n).
\end{equation}
It commutes with the magnitude differential, since isometries preserve
the conditions determining that differential.  On cochains the action is
$(g\varphi)(c)=\varphi(g^{-1}c)$.  Thus
$\MH_{n,\ell}(X;R)$ and $\MH^{n,\ell}(X;R)$ are $R\G$-modules, and the
latter form an algebra on which $\G$ acts by bigraded automorphisms.
\source This is the specialization of functoriality in
\cite{HW,LS,Hepworth}; the automorphism-group action is explicitly noted
in \cite[\S1.1]{HW}.  The formulas also give a direct verification over
every commutative $R$.

For the representation-valued Euler series we now take $R=\C$.
The complex representation ring $\RC(\G)$ is the Grothendieck group of
finite-dimensional complex representations, with
$[U\oplus V]=[U]+[V]$ and multiplication $[U][V]=[U\otimes_\C V]$.
Maschke's theorem identifies its additive group with the free abelian
group on the irreducible representations.  Its elements are
\defn{virtual representations}; their characters are defined by linear
extension of $\chi_V(g)=\tr(g\mid V)$.  Character theory and its
orthogonality formulas are used in their usual complex form
\cite[Chs.~1--2]{Serre}.

\begin{definition}\label{def:rep-mag}
 The representation-valued magnitude series is
 \begin{equation}\label{eq:rep-mag}
 \Mag_\G(X;q)=
 \sum_{\ell\in\Lambda}\sum_{n\geq0}
       (-1)^n[\MH_{n,\ell}(X;\C)]q^\ell
 \quad\in\RC(\G)[[q^\Lambda]].
 \end{equation}
\end{definition}
Each coefficient is a finite sum.  The same series results from using
cohomology: one may apply Euler--Poincar\'e to the dual permutation
cochain complex.  This assertion concerns the Euler series; no choice
of an identification between individual homology and cohomology spaces
is needed.

\source Definition~\ref{def:rep-mag} is the equivariant specialization
used in this note.  Its analogy with equivariant Ehrhart theory is with
the permutation representation on lattice points, whose character at
$g$ counts the lattice points fixed by $g$, as in
\cite[\S5, Lemma 5.2]{Stapledon}.  We do not attribute the present
magnitude construction to that reference.

\subsection{The character formula}

\begin{theorem}[Fixed-point character formula]\label{thm:character}
 For $g\in\G$, put $X^g=\{x\in X:gx=x\}$.  Then
 \begin{equation}\label{eq:character}
   \chi_{\Mag_\G(X;q)}(g)=\Mag(X^g;q).
 \end{equation}
\end{theorem}
\begin{proof}
Fix a length $\ell$.  Euler--Poincar\'e for a finite complex with a chain
endomorphism gives
\[
 \sum_n(-1)^n\tr(g\mid\MH_{n,\ell}(X;\C))
 =\sum_n(-1)^n\tr(g\mid C_{n,\ell}(X;\C)).
\]
The trace of a permutation matrix counts fixed basis elements.  By
\eqref{eq:tuple-action}, a tuple is fixed by $g$ exactly when every
entry lies in $X^g$.  Therefore the right side equals
$\sum_n(-1)^n|\mathcal T_{n,\ell}(X^g)|$, the coefficient of $q^\ell$
in ordinary magnitude.  This last identification is the ordinary Euler
characteristic theorem cited above.  The cochain argument gives the
same result, since $g$ and $g^{-1}$ fix the same tuples.
\end{proof}
\source This theorem is proved here from the standard permutation
trace identity and the ordinary magnitude Euler characteristic theorem.
It is not an asserted degree-by-degree identification of homology of
fixed spaces with fixed vectors in homology.

\begin{corollary}\label{cor:rep-specializations}
 The following statements hold over $\C$.
 \begin{enumerate}[label=(\roman*)]
 \item Applying dimension to $\Mag_\G(X;q)$ gives $\Mag(X;q)$.
 \item For $V\in\Irr(\G)$, the virtual multiplicity series of $V$ is
 \begin{equation}\label{eq:character-multiplicity}
  \sum_{n,\ell}(-1)^n
    \dim_\C\Hom_\G(V,\MH_{n,\ell}(X;\C))q^\ell
  =\frac1{|\G|}\sum_{g\in\G}\overline{\chi_V(g)}\Mag(X^g;q).
 \end{equation}
 \item In particular, the invariant Euler series is
 \begin{equation}\label{eq:averaged-mag}
  \sum_{n,\ell}(-1)^n\dim_\C\MH_{n,\ell}(X;\C)^\G q^\ell
  =\frac1{|\G|}\sum_{g\in\G}\Mag(X^g;q).
 \end{equation}
 \end{enumerate}
\end{corollary}
\begin{proof}
Evaluate \eqref{eq:character} at $1$, or take the character inner product
with $\chi_V$, respectively.  For (iii), take $V=\one$.
\end{proof}

The same arguments imply restriction compatibility for every subgroup
$H\leq\G$.  They also give
$\Mag_\G(X\times Y;q)=\Mag_\G(X;q)\Mag_\G(Y;q)$ for the diagonal
action and the sum metric.  Indeed, fixed points commute with that
product and ordinary magnitude is multiplicative: the magnitude
matrix of the product is the Kronecker product of the two magnitude
matrices.  Injectivity of complex characters then gives the identity.
These are deductions from Theorem~\ref{thm:character}.

\begin{example}[Two points interchanged]\label{ex:two-points}
Let $X=\{a,b\}$ with $d(a,b)=1$, and let $\G=C_2$ interchange them.
There are two proper alternating tuples in degree $n$, of length $n$,
and their differentials are zero.  Thus, over every commutative $R$,
\[
 \MH_{n,\ell}(X;R)=
 \begin{cases}R[C_2],&\ell=n,\\0,&\ell\ne n.\end{cases}
\]
Over $\C$, writing $\varepsilon$ for the sign representation,
\[
 \Mag_{C_2}(X;q)=\frac{[\one]+[\varepsilon]}{1+q}.
\]
Its character values are $2/(1+q)$ and $0$.  Its invariant Euler series
is $1/(1+q)$, whereas $X^{C_2}=\varnothing$ and $X/C_2$ is a singleton
of magnitude $1$.  The invariant chain $a+b$ explains why invariant
vectors must not be confused with fixed points.
\end{example}

\begin{remark}[Integral and modular information]
The integral theory is the bigraded $\Z\G$-module
$\MH_{*,*}(X;\Z)$, including its torsion, rather than just
\eqref{eq:rep-mag}.  Integral homology groups need not be free abelian,
so one cannot silently regard them as objects in a ring of integral
lattices.  Over a field of characteristic dividing $|\G|$, magnitude
homology still carries a group action, but semisimplicity and the above
complex character reconstruction are unavailable.  The series
\eqref{eq:rep-mag} is specifically a complex representation-ring
invariant of the integral construction.
\end{remark}

\section{Borel equivariant magnitude (co)homology}
\label{sec:borel}

\subsection{Definition over a commutative ring}

\coeff Throughout this subsection $R$ is any commutative unital ring;
in particular $R=\Z$, $\Q$, $\C$, or $\F_p$ is permitted.  Give $R$
the trivial $R\G$-module structure.  Derived tensor products and derived
Hom are taken over the indicated ring, not over the ordinary magnitude
homology groups.

\begin{definition}\label{def:borel}
 For each length $\ell$, define
 \begin{align}
 \MH^{\Bor,\G}_{n,\ell}(X;R)
   &=H_n\bigl(R\derivedtensor{R\G}C_{*,\ell}(X;R)\bigr),\label{eq:borel-hom}\\
 \MH^{n,\ell}_{\Bor,\G}(X;R)
   &=H^n\bigl(\RHom_{R\G}(R,C^{*,\ell}(X;R))\bigr).
      \label{eq:borel-cohom}
 \end{align}
\end{definition}
Here $H_n(\G;M)=\Tor_n^{R\G}(R,M)$ and
$H^n(\G;M)=\Ext_{R\G}^n(R,M)$ denote group homology and cohomology.
For example, a free group-bar resolution $Q_\bullet\to R$ gives the
models
\[
 \Tot(Q_\bullet\otimes_{R\G}C_{*,\ell}),\qquad
 \Tot\Hom_{R\G}(Q_\bullet,C^{*,\ell}),
\]
using a right resolution in the first formula and a left resolution in
the second.  The total differential has the usual tensor or Hom signs.
Group degree contributes to total (co)homological degree, and has
length zero.

\source These are applications of the standard derived coinvariant
and invariant constructions for group actions; see \cite[Chs.~I, III,
VII]{Brown}.  We use them here as definitions for magnitude complexes.
Their topological interpretation as reduced Borel constructions is
given in \S\ref{subsec:nerve}, and is consistent with the usual Borel
construction in \cite[\S2.1.4]{Guillou}.

The magnitude cochain product and a diagonal on the group-bar
resolution equip $\MH^{*,*}_{\Bor,\G}$ with an associative bigraded
product.  All signs use total cohomological degree, not the length
degree.  Its algebraic identification with the Yoneda product is
addressed in Theorem~\ref{thm:borel-skew}.

\begin{proposition}[Group (co)homology spectral sequences]
\label{prop:borel-ss}
 Over every commutative $R$, there are first-quadrant spectral sequences
 \begin{align}
 H_p(\G;\MH_{q,\ell}(X;R))
   &\Longrightarrow\MH^{\Bor,\G}_{p+q,\ell}(X;R),\label{eq:hom-ss}\\
 H^p(\G;\MH^{q,\ell}(X;R))
   &\Longrightarrow\MH^{p+q,\ell}_{\Bor,\G}(X;R).\label{eq:cohom-ss}
 \end{align}
 The cohomological sequence, summed over lengths, is multiplicative.
\end{proposition}
\begin{proof}
Filter the above double complexes by group degree.  The group-bar
terms are finite free $R\G$-modules because $\G$ is finite, so taking
the magnitude differential first gives the displayed $E_2$-pages.
There are finitely many terms in each total degree, and the magnitude
direction is bounded at every length.  The standard first-quadrant
convergence argument applies.  Bar diagonals respect the filtration,
which gives multiplicativity.  This is the usual hyperhomology
construction of \cite[Ch.~VII]{Brown}; over fields the corresponding
Ext-algebra statement is also \cite[Corollary 7.7]{Negron}.
\end{proof}

\subsection{Comparison with the representation-valued theory}

\begin{proposition}[When the group order is invertible]
\label{prop:averaging}
 If $|\G|\in R^\times$, then
 \begin{align}
 \MH^{\Bor,\G}_{n,\ell}(X;R)&\cong\MH_{n,\ell}(X;R)_\G,\\
 \MH^{n,\ell}_{\Bor,\G}(X;R)&\cong\MH^{n,\ell}(X;R)^\G.
 \end{align}
 The averaging map identifies coinvariants with invariants in this
 situation.  Over $\C$, the Borel Euler series equals
 \begin{equation}\label{eq:borel-euler}
 \frac1{|\G|}\sum_{g\in\G}\Mag(X^g;q).
 \end{equation}
\end{proposition}
\begin{proof}
The idempotent $|\G|^{-1}\sum_g g$ makes invariants an exact functor,
and the trivial $R\G$-module is projective on both sides.  Higher group
(co)homology vanishes, so Proposition~\ref{prop:borel-ss} collapses.
The isomorphism on coinvariants sends $[m]$ to
$|\G|^{-1}\sum_g gm$.  The Euler series then follows from
Corollary~\ref{cor:rep-specializations}.
\end{proof}

In (2.5), the subscript $\G$ on $\MH_{n,\ell}(X;R)_\G$ means
\defn{coinvariants}:
\[
 M_\G=M/\langle gm-m:g\in\G,\ m\in M\rangle.
\]
This is distinct from the superscript $M^\G$ of \defn{invariants},
the fixed submodule.  When $|\G|\in R^\times$, averaging identifies
the two, as stated in the proposition; over $\Z$ or in characteristic
dividing $|\G|$, they need not agree.

\begin{remark}[What specializes, and what does not]
For the trivial symmetry group both constructions recover ordinary
magnitude (co)homology.  For a nontrivial finite group in characteristic
zero, Borel cohomology recovers the \emph{invariant part} of the
representation-valued theory.  It does not recover the full
representation from untwisted coefficients alone.  Also,
\eqref{eq:borel-euler} need not equal either $\Mag(X^\G;q)$ or
$\Mag(X/\G;q)$; Example~\ref{ex:two-points} gives both failures.
\end{remark}

There is a precise recovery statement if one allows coefficients
twisted by irreducible representations.  For a finite-dimensional
complex representation $V$, put
\[
 \MH^{n,\ell}_{\Bor,\G}(X;V)
 :=H^n\RHom_{\C\G}(V,C^{*,\ell}(X;\C)).
\]
Equivalently this uses group cochains with values in
$V^*\otimes C^{*,\ell}$.  Semisimplicity gives
\begin{equation}\label{eq:twisted-recovery}
 \MH^{n,\ell}_{\Bor,\G}(X;V)
 =\Hom_\G(V,\MH^{n,\ell}(X;\C)),\qquad
 \MH^{n,\ell}(X;\C)
 \cong\bigoplus_{V\in\Irr(\G)}
       V\otimes\MH^{n,\ell}_{\Bor,\G}(X;V).
\end{equation}
The second is the evaluation isomorphism.  This deduction uses the
complex representation theory of \cite[Ch.~2]{Serre}; it is not a
modular reconstruction assertion.

\subsection{Stabilizers and modular phenomena}

\begin{proposition}[Length zero]\label{prop:length-zero}
 Over any commutative $R$,
 \begin{align}
 \MH^{n,0}_{\Bor,\G}(X;R)
   &\cong\bigoplus_{[x]\in X/\G}H^n(\G_x;R),\\
 \MH^{\Bor,\G}_{n,0}(X;R)
   &\cong\bigoplus_{[x]\in X/\G}H_n(\G_x;R).
 \end{align}
\end{proposition}
\begin{proof}
The length-zero magnitude complex is $R[X]$ in degree zero.  Decompose
it into the permutation modules $R[\G/\G_x]$.  For finite groups
induction and coinduction of the trivial module give this same
permutation module.  Apply the homological and cohomological forms of
Shapiro's lemma \cite[Ch.~III]{Brown}.
\end{proof}

For a single point, Borel magnitude (co)homology is group
(co)homology, all in length zero.  In particular, for $\G=C_p$,
\[
 H^n(C_p;\Z)=
 \begin{cases}
 \Z,&n=0,\\
 \Z/p,&n>0\text{ even},\\
 0,&n\text{ odd}.
 \end{cases}
\]
Over $\F_p$, $H^n(C_p;\F_p)$ is one-dimensional in every
nonnegative degree.  Both statements follow from the cyclic-group
resolution alternating $g-1$ and $1+g+\cdots+g^{p-1}$
\cite[Ch.~I, \S6]{Brown}.  Thus, in modular characteristic, even one
fixed point can produce unbounded cohomological degree at a fixed
length.  The finite alternating sum used in
Definition~\ref{def:rep-mag} is then unavailable; no Borel magnitude
Euler series is asserted in this generality.  Integral torsion is
likewise invisible to a rank-only Euler characteristic.

If the action on vertices is free, every proper tuple has trivial
stabilizer.  Its chain and cochain modules are sums of free permutation
$R\G$-modules.  These are acyclic for group homology and cohomology
(by Shapiro), so the Borel complexes can be computed by ordinary
coinvariant chains and invariant cochains even without inverting
$|\G|$.  This is a chain-level statement, not a claim that group
cohomology of each magnitude homology group vanishes.  These quotient
complexes still need not be the magnitude complex of the quotient
metric space, as Example~\ref{ex:two-points} shows.

\subsection*{\color{darkblue}Related literature}
Fu--Yau construct a Borel version of \emph{path} homology of directed
graphs \cite{FuYau}.  That is a related construction for a different
homology theory; none of the magnitude comparison results here is
being attributed to it.

\section{Distance algebras and skew group constructions}
\label{sec:skew}

\subsection{The algebra, the two modules, and grading conventions}

\coeff Unless a field hypothesis is explicitly added, $R$ is an
arbitrary commutative unital ring.  Let
\[
 A=\Sigma_X(R),\qquad J=A_{>0},\qquad S=A/J\cong\prod_{x\in X}R.
\]
The distance algebra and the Asao--Ivanov identifications are assumed.
We use $b_{xy}$ for its standard basis and $e_x=b_{xx}$.  The action
$g(b_{xy})=b_{gx,gy}$ preserves multiplication and grading, hence acts
on $A$, $J$, and $S$.  Write $\epsilon:A\to S$ for the quotient map.

\begin{definition}[Skew group algebra]
 The \defn{skew group algebra} is the $R$-module
 \[
  B=A\rtimes_R\G:=A\otimes_R R\G
 \]
 with multiplication determined by
 \[
  (a\otimes g)(b\otimes h)=a\,g(b)\otimes gh.
 \]
 Put $u_g=1\otimes g$, give each $u_g$ length zero, and set
 \begin{equation}\label{eq:T}
   T=B_0=S\rtimes\G=B/(J\rtimes\G).
 \end{equation}
\end{definition}
This is finite free of rank $|X|^2|\G|$ over $R$.  The module $S$
has the following \defn{separate} left and right $B$-structures:
\begin{equation}\label{eq:S-actions}
 (a\otimes g)\cdot s=\epsilon(a)g(s),\qquad
 s\cdot(a\otimes g)=g^{-1}(s\epsilon(a)).
\end{equation}
They are the structures used on the appropriate sides of Ext and Tor.
We do not assert that these two actions commute to give a $B$-bimodule.
For Ext we always use left modules.  The left module $S$ is the
permutation module, not the regular $T$-module.

For a graded module, $M[\ell]_d=M_{d-\ell}$, and
\[
 \Ext_A^{n,\ell}(M,N)
 =\Ext^n_{D\text{-}\mathrm{GrMod}(A)}(M,N[\ell]).
\]
Tor uses its internal homogeneous component of degree $\ell$.
The same convention applies to $B$.  In these conventions
\begin{equation}\label{eq:AI}
 \Tor^A_{n,\ell}(S,S)\cong\MH_{n,\ell}(X;R),\qquad
 E_R:=\Ext_A^{*,*}(S,S)\cong\MH^{*,*}(X;R).
\end{equation}
The last is an algebra isomorphism with the magnitude cup product.
\source These are \cite[Theorems 5.3 and 5.6]{AI}; their
Remark 5.7 explains the graded-opposite convention that would occur
with right-module Ext.  The equivariance in \eqref{eq:AI} follows
directly from their bar models, whose tuple actions are
\eqref{eq:tuple-action}.  This is the only extra observation needed
here.

\subsection{Borel theory as ordinary Tor and Ext over the skew algebra}

\begin{theorem}\label{thm:borel-skew}
 For every commutative $R$ there are natural identifications
 \begin{align}
  \MH^{n,\ell}_{\Bor,\G}(X;R)&\cong\Ext_B^{n,\ell}(S,S),\label{eq:borel-ext}\\
  \MH^{\Bor,\G}_{n,\ell}(X;R)&\cong\Tor^B_{n,\ell}(S,S).\label{eq:borel-tor}
 \end{align}
 The cohomological identification respects products.  Equivalently,
 lengthwise there are derived comparisons
 \begin{align}
  \RHom_B(S,S)&\simeq
   \RHom_{R\G}\bigl(R,\RHom_A(S,S)\bigr),\label{eq:derived-ext}\\
  S\derivedtensor{B} S&\simeq
   R\derivedtensor{R\G}\bigl(S\derivedtensor{A} S\bigr).
       \label{eq:derived-tor}
 \end{align}
 The inner derived objects carry their natural $\G$-actions.
\end{theorem}
\begin{proof}
Choose the equivariant left $A$-bar resolution $P_\bullet\to S$
from the left-module version of \cite[Remark 5.4, \S5.3]{AI}.  It is
$A$-projective and $R$-split as an augmented complex.  Let
$Q_\bullet\to R$ be a free left group-bar resolution.  The total complex
of $P_\bullet\otimes_R Q_\bullet$, with diagonal $\G$-action, is a
$B$-projective resolution of $S$.  To see projectivity, for any
equivariant $A$-module $P$ there is an isomorphism
\[
 B\otimes_A P\longrightarrow P\otimes_R R\G,
 \qquad (a u_g)\otimes p\longmapsto a\,g(p)\otimes g.
\]
If $P$ is $A$-projective, this is $B$-projective; each $Q_p$ is free
over $R\G$.  Exactness of the total augmentation follows from the
$R$-split bar augmentations.

Tensor--Hom adjunction identifies the Hom complex of this resolution
with
\[
 \Tot\Hom_{R\G}(Q_\bullet,\Hom_A(P_\bullet,S)).
\]
Using \eqref{eq:AI} at the level of bar cochains gives
\eqref{eq:derived-ext}.  Bar diagonals and the usual lifting
description of Yoneda composition identify its product with the
Borel cup product.  This is the bar-model mechanism of
\cite[Theorem 7.5 and Corollary 7.7]{Negron}, where the general
skew-group change-of-rings theorem is formulated over a field.  The argument here
works over $R$ because the specified bars are $R$-split and all the
induced terms are projective; no averaging or division is used.

For tensor products the degree-zero identity for equivariant right
and left modules is
$M\otimes_B N\cong(M\otimes_A N)_\G$, where
\[
 g\cdot(m\otimes n)=m u_{g^{-1}}\otimes u_g n.
\]
Resolving by the same induced projective terms gives
\eqref{eq:derived-tor}.  One can also verify acyclicity directly:
for a free right $B$-module $P$, the $\G$-module $P\otimes_A S$ is a
sum of induced modules from the trivial subgroup, since $S$ is finite
free over $R$.  It is therefore acyclic for coinvariants, and the
standard derived-composition argument applies.  The action above is
the diagonal tuple action after \eqref{eq:AI}.
\end{proof}

This theorem realizes the two spectral sequences of
Proposition~\ref{prop:borel-ss} as skew-group change-of-rings spectral
sequences.  It also identifies the length-zero groups in
Proposition~\ref{prop:length-zero} inside $\Ext_B(S,S)$ and
$\Tor^B(S,S)$.

\subsection{Recovery of the full equivariant theory}

\begin{proposition}[Induced coefficient module]\label{prop:mixed}
 For every commutative $R$,
 \begin{align}
 \Ext_B^{n,\ell}(T,S)&\cong\Ext_A^{n,\ell}(S,S),\label{eq:mixed-ext}\\
 \Tor^B_{n,\ell}(T,S)&\cong\Tor^A_{n,\ell}(S,S).
       \label{eq:mixed-tor}
 \end{align}
 Both recover the full $R\G$-module structure with the residual
 actions specified below.
\end{proposition}
\begin{proof}
There are canonical induction identifications
\[
 {}_B T\cong B\otimes_A S,\qquad T_B\cong S\otimes_A B.
\]
Since $B$ is free over $A$ on both sides, derived induction--restriction
adjunction proves the displayed formulas.  These are the usual
Eckmann--Shapiro comparisons, here proved simply by inducing a
projective resolution.

For Ext, right multiplication $r_g:t\mapsto t u_g$ on the left module
$T$ induces, by contravariance, a left $\G$-action on
$\Ext_B(T,S)$.  On $B\otimes_A P$ it lifts to
\[
 b\otimes p\longmapsto b u_g\otimes g^{-1}p.
\]
Under adjunction its pullback sends a cochain $f$ to
$p\mapsto g f(g^{-1}p)$, precisely the action in ordinary equivariant
cohomology.  For Tor, left multiplication by $u_g$ on the right module
$T$ lifts to $p\otimes b\mapsto gp\otimes u_g b$ on an induced
right resolution.  After tensoring with $S$, this acts diagonally on
both factors and gives the ordinary equivariant homology action.
\end{proof}

\begin{proposition}[Twisted coefficients]\label{prop:twisted-skew}
Let $R=\C$ and let $V$ be a finite-dimensional representation of $\G$.
Give $S\otimes_\C V$ the diagonal action, with $A$ acting on $S$.
Then
\[
 \Ext_B^{n,\ell}(S\otimes_\C V,S)
 \cong\Hom_\G\bigl(V,\MH^{n,\ell}(X;\C)\bigr).
\]
In particular, these Ext groups recover all multiplicity spaces of
the usual equivariant cohomology.
\end{proposition}
\begin{proof}
The derived comparison of Theorem~\ref{thm:borel-skew}, with first
argument $S\otimes V$, identifies the result with group cohomology
of $\Hom_\C(V,\RHom_A(S,S))$.  Exactness of complex invariants
gives the formula, which is \eqref{eq:twisted-recovery} in algebraic
form.
\end{proof}

\subsection{The Yoneda algebra and the role of simple modules}

\begin{theorem}[Yoneda skew-product identity]\label{thm:yoneda}
For the distance-algebra situation above, over every commutative $R$,
\begin{equation}\label{eq:yoneda-skew}
 \Ext_B^{*,*}(T,T)\cong E_R\rtimes\G
       \cong\MH^{*,*}(X;R)\rtimes\G
\end{equation}
as bigraded algebras.  Here $T$ is the module in \eqref{eq:T}, whether
or not it is semisimple.
\end{theorem}
\begin{proof}
Use an equivariant projective left $A$-bar resolution $P\to S$.
Then $B\otimes_A P\to T$ is a projective left $B$-resolution.
Equip $\End_A(P)$ with the conjugation $\G$-action.  There is an
isomorphism of differential graded algebras
\begin{equation}\label{eq:dg-skew}
 \End_A(P)\rtimes_R\G
   =\End_A(P)\otimes_R R\G
   \longrightarrow\End_B(B\otimes_A P),\qquad
 \Phi(f\otimes g)(b\otimes p)=b u_g^{-1}\otimes f(gp).
\end{equation}
Balancing follows from $a u_g^{-1}=u_g^{-1}g(a)$ and $A$-linearity
of $f$.  Decomposing $B$ into its finitely many group components shows
that every endomorphism on the right has a unique expression of this
form.  Multiplication is checked directly:
\[
 \Phi(f\otimes g)\Phi(h\otimes t)
      =\Phi\bigl(f\,(g\cdot h)\otimes gt\bigr).
\]
The group maps have degree zero and commute with the differentials,
so \eqref{eq:dg-skew} respects differential graded structures and
internal lengths.  Taking cohomology gives \eqref{eq:yoneda-skew}:
the underlying complex of a skew product is a finite direct sum of
copies of the original complex.  Projective resolutions identify
these endomorphism cohomologies with Yoneda Ext.  No semisimplicity
or division by $|\G|$ occurs in this proof.
\end{proof}
\source Over a field with $\operatorname{char}k\nmid|\G|$, the
standard general result is \cite[Theorem 10]{MV}.  The proof above
establishes the version actually used here over arbitrary $R$, using
the canonical equivariant bar resolution of $S$.  It should not be
read as a claim that $T$ is the semisimple quotient in modular
characteristic.

For example, the distinguished group elements retain the action by
\[
 (1\otimes g)(\alpha\otimes1)(1\otimes g^{-1})=(g\alpha)\otimes1.
\]
Thus recovery from \eqref{eq:yoneda-skew} uses its natural group
grading and distinguished copy of $R\G$, rather than just an
unmarked abstract algebra or its Morita class.

To discuss simples, take an algebraically closed field $k$.
The ideal $J\rtimes\G$ is nilpotent, so every simple $B$-module
comes from $T$.  Choosing a vertex representative in each orbit and
transporters to the other vertices gives the groupoid-algebra
decomposition
\begin{equation}\label{eq:groupoid}
 S\rtimes\G\cong
 \prod_{[x]\in X/\G}\operatorname{Mat}_{|\G x|}(k\G_x).
\end{equation}
Indeed, a basis element $e_yu_g$ has source $g^{-1}y$ and target $y$;
its multiplication is composition of such arrows.  Transporters
identify arrows within an orbit with matrix units labelled by
$\G_x$, proving \eqref{eq:groupoid}.  Hence simple modules are indexed
by pairs $([x],V)$ with $V$ an irreducible $k\G_x$-module.  The
permutation module $S$ selects the trivial representation of each
stabilizer.

If $\operatorname{char}k\nmid|\G|$, Maschke's theorem makes $T$
semisimple and $T=B/\rad B$.  It contains every simple $B$-module
with its regular-module multiplicity.  In this case
\eqref{eq:yoneda-skew} is the usual Yoneda algebra of the semisimple
quotient.  Passing to one copy of each simple gives a full idempotent
corner of this Ext algebra, not necessarily the identical algebra.
In modular characteristic $T$ can fail to be semisimple, although
Theorems~\ref{thm:borel-skew} and~\ref{thm:yoneda} remain valid with
the specified modules.

\subsection{Fixed-point spaces as corners}

For a subgroup $H\leq\G$, put
$e_H=\sum_{x\in X^H}e_x$ and $S_H=e_HS$.
\begin{proposition}\label{prop:fixed-corner}
For every commutative $R$,
\[
 \Sigma_{X^H}(R)\cong e_HAe_H,
 \qquad (e_HAe_H)_0\cong S_H.
\]
For $R=\C$ and $H=\langle g\rangle$, this gives the identity
\begin{equation}\label{eq:corner-character}
 \sum_{n,\ell}(-1)^n\tr\bigl(g\mid\Ext_A^{n,\ell}(S,S)\bigr)q^\ell
 =\sum_{n,\ell}(-1)^n
   \dim_\C\Ext_{e_HAe_H}^{n,\ell}(S_H,S_H)q^\ell.
\end{equation}
\end{proposition}
\begin{proof}
The corner has basis $b_{xy}$ with $x,y\in X^H$, and its multiplication
is exactly that for the restricted metric.  The Ext identity follows
from Theorem~\ref{thm:character} and \eqref{eq:AI} applied to this
metric space.  If there are no fixed vertices, both the corner
contribution and its magnitude groups are taken to be zero.
\end{proof}
\source The corner description is a direct deduction from the
distance-algebra multiplication of \cite[\S5.1]{AI};
\eqref{eq:corner-character} is proved here by the character formula.
Neither $e_HAe_H=A^H$ nor a general commutation of Ext with corners
is asserted.  The right side computes new Ext groups over the corner
algebra, and \eqref{eq:corner-character} is an Euler identity.

\begin{center}
\begin{tabular}{@{}p{0.38\textwidth}p{0.55\textwidth}@{}}
\toprule
Algebraic object & Interpretation and coefficients\\
\midrule
$\Ext_B(T,S)$ with residual action & Full usual equivariant
cohomology; every commutative $R$.\\[3pt]
$\Ext_B(S,S)$ & Borel cohomology; every commutative $R$.  It is
$E_R^\G$ when $|\G|\in R^\times$.\\[3pt]
$\Ext_B(T,T)$ & $E_R\rtimes\G$; every commutative $R$ with $T=B_0$.
It is the semisimple-quotient Yoneda algebra over nonmodular fields.\\[3pt]
$\Ext_{e_HAe_H}(S_H,S_H)$ & Cohomology of the restricted fixed-point
metric space; every commutative $R$.\\
\bottomrule
\end{tabular}
\end{center}

\section{An integral Burnside-ring analogue}
\label{sec:burnside}

\subsection{Finite group sets, orbits, and stabilizers}

This section introduces the necessary Burnside-ring terminology from
the beginning.  The constructions in
\S\S\ref{subsec:burnside-ring}--\ref{subsec:burnside-series} are
integral: they use finite sets and integers and have no ground-field
choice.  Standard references for the background are
\cite[\S\S2--3]{Bouc}; the metric-space applications are proved here.

A \defn{finite $\G$-set} is a finite set $U$ with a left action of
$\G$.  A map $f:U\to V$ is \defn{equivariant} if $f(gu)=gf(u)$.
An isomorphism of $\G$-sets is an equivariant bijection.  For $u\in U$,
its \defn{orbit} is $\G u=\{gu:g\in\G\}$, and its
\defn{stabilizer} is $\G_u=\{g:gu=u\}$.  For a subgroup $H\leq\G$,
\[
 U^H=\{u:hu=u\text{ for every }h\in H\}
\]
is its set of \defn{fixed points}.  A set is \defn{transitive} if it consists
of one orbit.  Every transitive $\G$-set is isomorphic to $\G/H$ with
the left coset action, for a stabilizer $H$, via $gH\mapsto gu$.
Two such sets $\G/H$ and $\G/K$ are isomorphic exactly when $H$ and
$K$ are conjugate.  Consequently any finite $\G$-set decomposes
uniquely, up to order and isomorphism, into transitive orbits.

Write $\mathscr S(\G)$ for a chosen set of representatives of the
conjugacy classes of subgroups.  A subgroup $N\leq\G$ is
\defn{normal}, written $N\triangleleft\G$, if $gNg^{-1}=N$ for every
$g\in\G$.  Write
$N_\G(H)=\{g:gHg^{-1}=H\}$ for the \defn{normalizer}, and
$W_\G(H)=N_\G(H)/H$ for the \defn{Weyl group}.  These definitions
apply to all subgroups, including $1$ and $\G$.

\subsection{The Burnside ring and its structure maps}
\label{subsec:burnside-ring}

\begin{definition}[Burnside ring]
The \defn{Burnside ring} $\Burn(\G)$ is the abelian group generated by
symbols $[U]$ for isomorphism classes of finite $\G$-sets, subject to
\[
 [U\sqcup V]=[U]+[V].
\]
Multiplication is $[U][V]=[U\times V]$, using the diagonal action.
The zero is $[\varnothing]$ and the unit is $[\G/\G]$, a single fixed
point.  Arbitrary elements, including differences of sets, are called
\defn{virtual finite $\G$-sets}.
\end{definition}
Distributivity and commutativity follow from the corresponding
equivariant bijections of finite sets.  Orbit decomposition gives
\begin{equation}\label{eq:orbit-basis}
 \Burn(\G)=\bigoplus_{H\in\mathscr S(\G)}\Z[\G/H].
\end{equation}
In particular, this is a free abelian group of finite rank.  An
element is represented by an actual finite set precisely when all
its orbit coefficients are nonnegative.  Such elements are sometimes
called \defn{effective}.

The product in this basis is described by double cosets:
\begin{equation}\label{eq:burnside-product}
 [\G/H][\G/K]
  =\sum_{g\in H\backslash\G/K}[\G/(H\cap gKg^{-1})].
\end{equation}
Here $H\backslash\G/K$ means the set of subsets $HgK$, and one
representative $g$ is chosen from each.  To prove the formula, move
the first coordinate of a pair of cosets to $H$.  Its remaining
orbits are represented by $(H,gK)$, whose stabilizer is
$H\cap gKg^{-1}$.  This is the standard orbit-product calculation
of \cite[\S\S2.2--3.1]{Bouc}.

The following maps will be useful.
\begin{itemize}
\item The \defn{cardinality augmentation}
$\operatorname{aug}:\Burn(\G)\to\Z$ sends $[U]$ to $|U|$.
It is a ring homomorphism.
\item \defn{Restriction} $\Res_H^\G:\Burn(\G)\to\Burn(H)$ forgets
part of the action and is a ring homomorphism.
\item \defn{Induction} $\Ind_H^\G:\Burn(H)\to\Burn(\G)$ sends
$[V]$ to $[\G\times_H V]$, where
$(gh,v)\sim(g,hv)$.  It is additive but generally not multiplicative.
For $K\leq H$, it sends $[H/K]$ to $[\G/K]$.
\item \defn{Conjugation} by $g$ transports an $H$-action to a
$gHg^{-1}$-action: $ghg^{-1}$ acts as $h$ did.  It gives an isomorphism
$c_g:\Burn(H)\to\Burn(gHg^{-1})$.
\end{itemize}
These maps satisfy transitivity, conjugation compatibility, the
double-coset formula, and Frobenius reciprocity.  For clarity, the
last two identities are
\begin{align}
 \Res_K^\G\Ind_H^\G
 &=\sum_{g\in K\backslash\G/H}
   \Ind_{K\cap gHg^{-1}}^K\,c_g\,
   \Res_{g^{-1}Kg\cap H}^{H},\label{eq:mackey}\\
 a\,\Ind_H^\G(b)&=\Ind_H^\G(\Res_H^\G(a)b).
       \label{eq:frobenius}
\end{align}
They follow by decomposing the induced sets into $K$-orbits and by
the bijection $(u,[g,v])\mapsto[g,(g^{-1}u,v)]$, respectively.
The compatible additive restriction, induction and conjugation data
are called a \defn{Mackey functor}; in the presence of ring structures
and \eqref{eq:frobenius} they give the Burnside \defn{Green functor}.
We use these terms only for this concrete system of maps.  The
general definitions and these identities are in
\cite[\S2.2, Proposition 3.1.3, \S5]{Bouc}.

\subsection{Marks and the table of marks}

\begin{definition}[Marks]
For $H\leq\G$, the $H$-\defn{mark} is the ring homomorphism
\[
 m_H:\Burn(\G)\longrightarrow\Z,
 \qquad m_H([U])=|U^H|.
\]
It depends only on the conjugacy class of $H$.  The combined
\defn{mark homomorphism} is
\begin{equation}\label{eq:marks}
 m:\Burn(\G)\longrightarrow\prod_{H\in\mathscr S(\G)}\Z,
 \qquad a\longmapsto(m_H(a))_H.
\end{equation}
The target is also called the \defn{ghost ring}, with coordinatewise
addition and multiplication.
\end{definition}
The homomorphism property follows because fixed points commute with
disjoint union and Cartesian product.  Note that $m_1$ is cardinality,
whereas $m_\G$ counts points fixed by the whole group.

\begin{proposition}[Injectivity and rational reconstruction]
\label{prop:marks}
The map $m$ is injective and has finite cokernel.  After tensoring
with $\Q$, it is an isomorphism
\[
 \Q\otimes_\Z\Burn(\G)\cong
       \prod_{H\in\mathscr S(\G)}\Q.
\]
\end{proposition}
\begin{proof}
In the orbit basis the matrix is the \defn{table of marks}
\begin{equation}\label{eq:table-marks}
 M_{H,K}=m_H([\G/K])
 =|\{gK:g^{-1}Hg\leq K\}|.
\end{equation}
An entry can be nonzero only if $H$ is conjugate to a subgroup of $K$.
Order the subgroup representatives by nonincreasing order.  The
matrix is triangular, with diagonal entry
$M_{H,H}=|N_\G(H):H|=|W_\G(H)|$, a positive integer.  It has
nonzero determinant, which proves all assertions.
\end{proof}
\source This is the classical mark theorem, with the standard
triangular-matrix proof; see \cite[\S3.2]{Bouc}.  In particular,
arbitrary integer lists of fixed-point counts need not come from a
virtual finite set: they must belong to the integral image of this
matrix.

\begin{example}[A cyclic group of prime order]\label{ex:cp-marks}
For $\G=C_p$, put $1=[\G/\G]$ and $t=[\G/1]$.
Then
\[
 \Burn(C_p)=\Z\,1\oplus\Z\,t,\qquad t^2=pt,
\]
and
\[
 m_\G(a+bt)=a,\qquad m_1(a+bt)=a+pb.
\]
Thus an integer pair $(z_1,z_\G)$ is a mark vector precisely when
$z_1\equiv z_\G\pmod p$.  This is a direct calculation from the
two orbit types, not an additional assumption on marks.
\end{example}

\subsection{Linearization and information lost in characters}

\begin{definition}[Complex linearization]
The ring homomorphism
\[
 L_\C:\Burn(\G)\longrightarrow\RC(\G),\qquad
 [U]\longmapsto[\C[U]]
\]
is called \defn{linearization}.  Its character satisfies
\begin{equation}\label{eq:linearization-marks}
 \chi_{L_\C(a)}(g)=m_{\langle g\rangle}(a).
\end{equation}
\end{definition}
Indeed, the trace counts elements of $U$ fixed by $g$.  Thus complex
characters see exactly the marks at cyclic subgroups.  In
particular, $a\in\ker L_\C$ if and only if all its cyclic-subgroup
marks vanish.  This is the usual permutation-character construction
\cite[Ch.~1]{Serre}, combined with the definition of marks.  It
does not imply that every complex representation is a linearized
finite set.

\begin{example}[A nonzero element killed by linearization]
Let $\G=C_2\times C_2$, with its three order-two subgroups
$H_1,H_2,H_3$.  The element
\[
 a=[\G/1]-\sum_{i=1}^3[\G/H_i]+2[\G/\G]
\]
has $m_1(a)=4-6+2=0$ and $m_{H_i}(a)=0-2+2=0$ for every $i$.
These are all cyclic subgroups, so $L_\C(a)=0$.  However,
$m_\G(a)=2$, so $a\ne0$.  This calculation shows explicitly how
whole-subgroup fixed-point information can be lost upon passing to
complex representations.
\end{example}

For every commutative $R$ the operation $U\mapsto R[U]$ is still
available on finite sets.  The comparison with ordinary complex
characters in \eqref{eq:linearization-marks} specifically uses
$R=\C$.  We make no identification of $\Burn(\G)$ with a ring of
$R\G$-modules: linearization can lose information even over $\C$.

\subsection{The Burnside-valued magnitude series}
\label{subsec:burnside-series}

\begin{definition}\label{def:burnside-mag}
The integral Burnside-valued magnitude series is
\begin{equation}\label{eq:burnside-mag}
 \Mag^{\Burn}_\G(X;q)
  =\sum_{\ell\in\Lambda}\sum_{n\geq0}
           (-1)^n[\mathcal T_{n,\ell}(X)]q^\ell
  \quad\in\Burn(\G)[[q^\Lambda]].
\end{equation}
\end{definition}
The brackets here mean finite $\G$-sets, not modules or homology
groups.  The boundedness established in \S1 makes every coefficient
well-defined integrally.

\begin{theorem}[Marks and linearization of magnitude]
\label{thm:burnside-mag}
For every subgroup $H\leq\G$,
\begin{equation}\label{eq:burnside-fixed}
 m_H\bigl(\Mag^{\Burn}_\G(X;q)\bigr)=\Mag(X^H;q).
\end{equation}
Moreover,
\begin{equation}\label{eq:burnside-linear}
 L_\C\bigl(\Mag^{\Burn}_\G(X;q)\bigr)=\Mag_\G(X;q).
\end{equation}
Thus the collection of ordinary magnitudes of all restricted
fixed-point spaces uniquely determines the Burnside-valued series.
\end{theorem}
\begin{proof}
Coordinatewise fixedness gives
$\mathcal T_{n,\ell}(X)^H=\mathcal T_{n,\ell}(X^H)$.
Applying a mark to \eqref{eq:burnside-mag} therefore gives the
ordinary chain Euler series, proving \eqref{eq:burnside-fixed}.
Applying linearization gives the virtual Euler characteristic of
the complex permutation chain modules.  Euler--Poincar\'e in
$\RC(\G)$ gives \eqref{eq:burnside-linear}.  Uniqueness follows
coefficientwise from Proposition~\ref{prop:marks}.
\end{proof}
\source This theorem is proved here by applying the standard
Burnside-ring constructions to the magnitude tuple sets.  The only
magnitude input is the ordinary Euler characteristic theorem
\cite{HW,LS}.  In particular, existence is established integrally
by \eqref{eq:burnside-mag}; one does not merely postulate a rational
mark vector and assume it has an integral lift.

Several consequences follow immediately.
\begin{enumerate}[label=(\roman*)]
\item Cardinality recovers $\Mag(X;q)$, while the whole-group mark
recovers $\Mag(X^\G;q)$.
\item Restriction is compatible with restricting the group action.
The sum-metric product satisfies
\[
 \Mag^{\Burn}_\G(X\times Y;q)
  =\Mag^{\Burn}_\G(X;q)\Mag^{\Burn}_\G(Y;q).
\]
For the product assertion apply every mark, use multiplicativity of
ordinary magnitude, and use injectivity of marks.
\item For $\G=C_p$, Example~\ref{ex:cp-marks} gives the coefficientwise
integral congruence
\begin{equation}\label{eq:cp-congruence}
 \Mag(X;q)\equiv\Mag(X^{C_p};q)\pmod p.
\end{equation}
Indeed,
\[
 \Mag^{\Burn}_{C_p}(X;q)
 =\Mag(X^{C_p};q)\,1+
   \frac{\Mag(X;q)-\Mag(X^{C_p};q)}p\,t,
\]
and the quotient has integral coefficients.
\item If $\G$ acts freely on $X$, every tuple set is free, so
\[
 \Mag^{\Burn}_\G(X;q)=
       [\G/1]\,\frac{\Mag(X;q)}{|\G|},
\]
where the quotient is again coefficientwise integral.  This does not
identify it with magnitude of the quotient metric space.
\end{enumerate}
All four are deductions from the construction and mark theorem.
The noninjectivity of linearization shows that the Burnside target
can retain more information than the representation-ring target;
it does not by itself supply two metric examples distinguished by
their Burnside magnitude but not their representation-valued magnitude.

\subsection{An equivariant topological model}
\label{subsec:nerve}

To explain a possible homological refinement, we use the standard
pointed simplicial magnitude model, solely to expose its fixed-point
structure.  For graphs this is \cite[Definition 40 and Lemma 42]{HW};
the same description for finite metric spaces follows verbatim from
the length condition and is compatible with \cite{LS}.

For a length $\ell$, let $N_\ell(X)$ have a basepoint in each
simplicial degree and, in degree $n$, the tuples of total length
$\ell$, now allowing consecutive repetitions.  A face deletes an
entry if length is preserved and otherwise gives the basepoint;
a degeneracy repeats an entry.  Put
$K_\ell(X)=|N_\ell(X)|$, its pointed geometric realization.
Its non-basepoint nondegenerate $n$-simplices are precisely
$\mathcal T_{n,\ell}(X)$, and its reduced normalized chain complex is
$C_{*,\ell}(X;\Z)$.

A \defn{$\G$-CW complex} is built by attaching cells
$\G/H\times D^n$, with $\G$ acting on the coset factor and trivially
on the disk, along their boundaries.  In this model, the stabilizer
of a tuple acts trivially on its simplex: it fixes every vertex
position.  Thus $K_\ell(X)$ is a finite pointed $\G$-CW complex.
It is finite because only finitely many nondegenerate tuples have
the given length.

\begin{proposition}\label{prop:nerve-fixed}
For every $H\leq\G$,
\[
 N_\ell(X)^H=N_\ell(X^H),\qquad
 K_\ell(X)^H\cong K_\ell(X^H).
\]
Consequently, for every commutative $R$,
$\widetilde H_n(K_\ell(X)^H;R)\cong\MH_{n,\ell}(X^H;R)$.
\end{proposition}
\begin{proof}
The simplicial identity is coordinatewise fixedness, including
degeneracies and the basepoint.  In a geometric realization every
point is represented in the interior of a unique nondegenerate
simplex.  If it is fixed by $H$, this simplex is stabilized by $H$;
in the present model the stabilizer fixes it pointwise.  This proves
the realization identity.  Reduced cellular chains give the final
claim.
\end{proof}

For a finite pointed $\G$-CW complex $K$, let $\mathscr E_n$ be its
finite $\G$-set of non-basepoint $n$-cells.  Its \defn{reduced
equivariant Euler characteristic} is
\[
 \widetilde\chi^\G(K)=\sum_n(-1)^n[\mathscr E_n]\in\Burn(\G).
\]
Every mark is $\widetilde\chi(K^H)$, so injectivity of marks shows
that this expression is independent of the chosen finite
$\G$-CW structure and invariant under equivariant homotopy equivalence.
This is the usual equivariant Euler construction, in the form of
\cite[\S4]{Bouc}, with the basepoint subtracted.  Here it gives
\begin{equation}\label{eq:burnside-topological}
 \Mag^{\Burn}_\G(X;q)
   =\sum_\ell\widetilde\chi^\G(K_\ell(X))q^\ell.
\end{equation}

For comparison with \S2, choose a contractible free $\G$-CW complex
$E\G$ and put $B\G=E\G/\G$.  The usual reduced Borel model is the
pair
\[
 \bigl(E\G\times_\G K_\ell(X),\ E\G\times_\G\{*\}\bigr).
\]
Its relative chains are modeled by
$C_*(E\G;R)\otimes_{R\G}\widetilde C_*(K_\ell(X);R)$, using
cellular or simplicial Eilenberg--Zilber.  They compute
\eqref{eq:borel-hom}; relative cochains compute
\eqref{eq:borel-cohom}.  The basepoint subspace is $B\G$ and must be
removed in this relative/reduced construction.  For $X$ a point,
$K_0(X)=S^0$, so this still yields the ordinary (unreduced) group
(co)homology of $B\G$, as in \S2.  This explains the Borel terminology
and the reduction convention; see \cite[\S2.1.4]{Guillou} for the
underlying topological construction.

\subsection{A possible Bredon refinement: definitions and scope}
\label{subsec:bredon}

The pointed complexes above allow one to apply Bredon (co)homology.
We describe the necessary terminology explicitly.  The standard
construction is in \cite[\S2.1.6]{Guillou}; its application to
$K_\ell(X)$ is the proposal made here.

\begin{definition}[Orbit category and coefficient systems]
The \defn{orbit category} $\Orb_\G$ has objects $\G/H$, one for
each subgroup, and morphisms the equivariant maps between these
sets.  A \defn{contravariant coefficient system over $R$} is a
functor $\Orb_\G^{\mathrm{op}}\to R\text{-}\mathrm{Mod}$.
A \defn{covariant coefficient system} is a functor
$\Orb_\G\to R\text{-}\mathrm{Mod}$.
\end{definition}
A map $f_g:\G/H\to\G/K$ is determined by $f_g(H)=gK$, and exists
precisely when $g^{-1}Hg\leq K$.  It induces a map
$K_\ell(X)^K\to K_\ell(X)^H$, $z\mapsto gz$.  Thus
\begin{equation}\label{eq:orbit-chain}
 \underline C_{n,\ell}(X;R)(\G/H)
  =\widetilde C_n(K_\ell(X)^H;R)
  =R[\mathcal T_{n,\ell}(X^H)]
\end{equation}
is a chain complex of contravariant coefficient systems.  The
chains here are generated by fixed cells, not invariant linear
combinations of all cells.  A cell orbit $\G/K$ contributes the
representable system $R[\Hom_{\Orb_\G}(-,\G/K)]$.  Hence this is
precisely the usual reduced Bredon cellular chain system.  Its
objectwise homology is the diagram
$\G/H\mapsto\MH_{n,\ell}(X^H;R)$.

For a covariant system $M$, define the tensor product over the orbit
category by
\begin{equation}\label{eq:coend}
 \underline C_{n,\ell}\otimes_{\Orb_\G}M
 =\left(\bigoplus_H
       \underline C_{n,\ell}(\G/H)\otimes_R M(\G/H)\right)
   \big/\left\langle
      \underline C(f)c\otimes m-c\otimes M(f)m
     \right\rangle.
\end{equation}
Here $f:\G/H\to\G/K$, $c\in\underline C_{n,\ell}(\G/K)$,
and $m\in M(\G/H)$.  This quotient is often called a \defn{coend};
the displayed relations are its full meaning here.  For a
contravariant system $N$, let
$\Nat_R(\underline C_{n,\ell},N)$ denote the $R$-module of families
of $R$-linear maps commuting with all orbit-category morphisms.

\begin{definition}[Bredon magnitude groups]
For any commutative $R$ and coefficient systems as above, put
\begin{align}
 \MH^{\Br,\G}_{n,\ell}(X;M)
   &=H_n(\underline C_{*,\ell}(X;R)\otimes_{\Orb_\G}M),\\
 \MH^{n,\ell}_{\Br,\G}(X;N)
   &=H^n(\Nat_R(\underline C_{*,\ell}(X;R),N)).
\end{align}
\end{definition}
These are ordinary reduced Bredon groups of $K_\ell(X)$, with
length retained.  When $\G=1$ and the coefficient system is $R$,
they reduce to ordinary magnitude (co)homology.  Their equivariant
homotopy invariance as invariants of $K_\ell(X)$ is the usual
Bredon invariance, not a new assertion of metric homotopy invariance.

One natural integral choice uses the Burnside coefficient systems:
at $\G/H$ take $\Burn(H)$.  On
$f_g:\G/H\to\G/K$, the covariant structure map is
\[
 \Burn(H)\xrightarrow{c_{g^{-1}}}\Burn(g^{-1}Hg)
       \xrightarrow{\Ind_{g^{-1}Hg}^{K}}\Burn(K),
\]
and the contravariant structure map is
\[
 \Burn(K)\xrightarrow{\Res_{g^{-1}Hg}^{K}}\Burn(g^{-1}Hg)
       \xrightarrow{c_g}\Burn(H).
\]
These are well-defined on the coset $gK$ and respect composition;
inner conjugations act trivially on isomorphism classes of group
sets.  They are the covariant and contravariant parts of the
Burnside Mackey functor introduced above.  Over a general commutative
$R$, replace them by $R\otimes_\Z\Burn(H)$.

\subsection*{\color{darkblue}What is established and what remains a refinement problem}
The integral series \eqref{eq:burnside-mag}, its mark formula, its
linearization, and its realization as \eqref{eq:burnside-topological}
are established by the proofs in this section.  The Bredon groups
are also well-defined by the displayed construction.  However, a
Bredon homology group is an abelian group (or part of a coefficient
system), not automatically an element of $\Burn(\G)$.  We have not
identified an alternating sum of their classes in some unspecified
Grothendieck group with \eqref{eq:burnside-mag}.  Nor does the
ordinary $\Z\G$-module $\MH_{n,\ell}(X;\Z)$ have a canonical
Burnside class: homology need not possess a permutation basis, and
linearization is not injective.

The orbit-category construction retains all subgroup fixed-point
complexes and the maps between them.  The skew group algebra in
\S3 organizes equivariant $A$-modules and the Borel/representation
comparisons, while the individual fixed-point distance algebras
$e_HAe_H$ realize the values at subgroups.  A further algebraic
description of the complete orbit-category construction, including
its restriction and induction structures, is a reasonable next
problem.  No identification of that entire structure with ordinary
Tor or Ext over the single skew group algebra is claimed here.

\begin{thebibliography}{AI24}
\bibitem[AI24]{AI}
Y.~Asao and S.~O.~Ivanov,
\emph{Magnitude homology is a derived functor},
arXiv:2402.14466v3 (2024).
\href{https://arxiv.org/html/2402.14466v3}{Full text}.
Theorem and section numbers in this note refer to version 3.

\bibitem[Bo00]{Bouc}
S.~Bouc, \emph{Burnside rings}, in \emph{Handbook of Algebra},
vol.~2, Elsevier, 2000, pp.~739--804.
\href{https://doi.org/10.1016/S1570-7954(00)80043-1}{DOI}.
\href{https://lamfa.u-picardie.fr/bouc/burnamscorr.pdf}{Author's corrected text (2016)}.
Numbered references are to this corrected text.

\bibitem[Br82]{Brown}
K.~S.~Brown, \emph{Cohomology of Groups}, Graduate Texts in
Mathematics 87, Springer, 1982.
\href{https://doi.org/10.1007/978-1-4684-9327-6}{DOI}.

\bibitem[FY26]{FuYau}
X.~Fu and S.-T.~Yau,
\emph{Szczarba's twisted shuffle and equivariant path homology of
directed graphs}, arXiv:2603.09406 (2026).
\href{https://arxiv.org/abs/2603.09406}{Preprint}.

\bibitem[Gu20]{Guillou}
B.~Guillou, \emph{Class Notes: Equivariant Homotopy and Cohomology,
Math 751}, University of Kentucky,
Fall 2020, especially \S\S2.1.4--2.1.6.
\href{https://www.ms.uky.edu/~guillou/F20/751Notes.pdf}{Lecture notes}.

\bibitem[He18]{Hepworth}
R.~Hepworth, \emph{Magnitude cohomology},
arXiv:1807.06832.
\href{https://arxiv.org/abs/1807.06832}{Preprint and publication record}.

\bibitem[HW17]{HW}
R.~Hepworth and S.~Willerton,
\emph{Categorifying the magnitude of a graph},
Homology, Homotopy and Applications \textbf{19} (2017), no.~2, 31--60.
\href{https://doi.org/10.4310/HHA.2017.v19.n2.a3}{DOI};
\href{https://arxiv.org/pdf/1505.04125v2}{arXiv:1505.04125v2}.

\bibitem[LS21]{LS}
T.~Leinster and M.~Shulman,
\emph{Magnitude homology of enriched categories and metric spaces},
Algebraic \& Geometric Topology \textbf{21} (2021), 2175--2221.
\href{https://doi.org/10.2140/agt.2021.21.2175}{DOI};
\href{https://arxiv.org/abs/1711.00802v4}{arXiv:1711.00802v4}.
Numbering follows version 4.

\bibitem[MV01]{MV}
R.~Mart\'inez-Villa,
\emph{Skew group algebras and their Yoneda algebras},
Mathematical Journal of Okayama University \textbf{43} (2001), 1--16.
\href{https://www.math.okayama-u.ac.jp/mjou/mjou43/_Martinez-Villa.pdf}{Full text}.

\bibitem[Ne14]{Negron}
C.~Negron,
\emph{Spectral sequences for the cohomology rings of a smash product},
arXiv:1401.3551v2 (2014).
\href{https://arxiv.org/html/1401.3551v2}{Full text (version 2 numbering)}.

\bibitem[Se77]{Serre}
J.-P.~Serre, \emph{Linear Representations of Finite Groups},
Graduate Texts in Mathematics 42, Springer, 1977.
\href{https://doi.org/10.1007/978-1-4684-9458-7}{DOI}.

\bibitem[St11]{Stapledon}
A.~Stapledon, \emph{Equivariant Ehrhart theory},
Advances in Mathematics \textbf{226} (2011), no.~4, 3622--3654.
\href{https://doi.org/10.1016/j.aim.2010.10.019}{DOI};
\href{https://arxiv.org/html/1003.5875v3}{arXiv:1003.5875v3}.
\end{thebibliography}
\end{document}
